\section{思想引领}

\subsection{党建活动}

此处介绍支部党员构成情况（正式党员、预备党员、入党积极分子人数），
以及班内党建氛围和同学们参与党建活动的积极性。

此处介绍本学年班级同学自发组织的党建学习活动，
如在党支部组织生活会上分享学习心得等。

% \begin{figure}[H]
%     \centering
%     \includegraphics[width=0.6\textwidth]{assets/example.jpg}
%     \caption{XX 同学分享学习心得}
% \end{figure}

此处介绍秋季学期主题党团日中的微党课环节，包括主题、形式和内容。

此处介绍春季学期的主题党团日活动，如前往红色教育基地参观、重温入党誓词等。

\subsection{主题党团日}

\subsubsection{大三上学期主题党团日}

此处描述活动的时间、地点、主题和联合举办情况（如联合其他团支部），
以及活动的核心内容（如生涯发展规划与外推经验分享）。

此处列举邀请到的嘉宾及其身份背景。

此处描述活动的开场环节，如主持人介绍嘉宾、辅导员进行学风和安全纪律宣讲等。

% \begin{figure}[H]
%     \centering
%     \includegraphics[width=0.45\textwidth]{assets/example.jpg}
%     \includegraphics[width=0.45\textwidth]{assets/example.jpg}
%     \caption{大三上学期主题党团日活动现场。左：XX 同学主持；右：XX 辅导员宣讲}
% \end{figure}

此处描述微党课环节的内容和形式。

% \begin{figure}[H]
%     \centering
%     \includegraphics[width=0.6\textwidth]{assets/example.jpg}
%     \caption{XX 同学分享微党课}
% \end{figure}

此处描述嘉宾分享环节的内容，如研究方向介绍、保研流程讲解、前沿工作分享等。

% \begin{figure}[H]
%     \centering
%     \includegraphics[width=0.45\textwidth]{assets/example.jpg}
%     \includegraphics[width=0.45\textwidth]{assets/example.jpg}
%     \caption{主题党团日嘉宾分享。左：XX 研究员；右：XX 教授}
% \end{figure}

此处描述自由提问环节和老师的总结点评。

% \begin{figure}[H]
%     \centering
%     \includegraphics[width=0.8\textwidth]{assets/example.jpg}
%     \caption{大三上学期主题党团日合影}
% \end{figure}

此处总结本次活动的整体效果和同学们的收获。

\subsubsection{大三下学期主题党团日}

此处描述活动的时间、地点、主题和形式（如实地探访、户外锻炼等）。

此处描述同学们集合出发的情况。

此处描述活动的主要内容，如参观纪念广场、聆听讲解等。

此处描述入党誓词重温环节。

% \begin{figure}[H]
%     \centering
%     \includegraphics[width=0.45\textwidth]{assets/example.jpg}
%     \includegraphics[width=0.45\textwidth]{assets/example.jpg}
%     \caption{大三下学期主题党团日活动现场}
% \end{figure}

此处描述自由参观环节。

% \begin{figure}[H]
%     \centering
%     \includegraphics[width=0.8\textwidth]{assets/example.jpg}
%     \caption{主题党团日活动}
% \end{figure}

此处描述户外活动环节（如登山等）和活动的整体总结。

% \begin{figure}[H]
%     \centering
%     \includegraphics[width=0.8\textwidth]{assets/example.jpg}
%     \caption{大三下学期主题党团日合影}
% \end{figure}

\subsection{志愿、实践、社工}

\subsubsection{志愿}

此处开篇介绍班级同学积极参与志愿活动的整体情况。

\begin{table}[H]
    \centering
    \zihao{-3}
    \begin{tabular}{||c|c||}
        \hline
        姓名   & 志愿者分级   \\
        \hline
        XXX & 五星级志愿者 \\
        XXX & 五星级志愿者 \\
        XXX & 四星级志愿者 \\
        XXX & 一星级志愿者 \\
        \hline
    \end{tabular}
    \caption{XX 班同学志愿者分级}
\end{table}

\begin{table}[H]
    \centering
    \zihao{-3}
    \begin{tabular}{||c|c||}
        \hline
        姓名   & 志愿时长   \\
        \hline
        XXX & XX 小时 \\
        XXX & XX 小时 \\
        XXX & XX 小时 \\
        \hline
    \end{tabular}
    \caption{XX 班同学志愿服务时长}
\end{table}

\subsubsection{实践}

此处介绍班级同学参加寒暑假社会实践的整体情况，包括参与人次和支队长情况。

\begin{table}[H]
    \centering
    \zihao{-3}
    \begin{tabular}{||c|c|c||}
        \hline
        \hline
        姓名   & 支队           & 奖项 \\
        \hline
        XXX & XX 实践支队     & 金奖 \\
        \hline
        XXX & XX 海外实践支队 & 金奖 \\
        \hline
        XXX & XX 实践支队     & 银奖 \\
            & XX 实践支队     &      \\
        \hline
        XXX & XX 实践支队     & 铜奖 \\
        \hline
        \hline
    \end{tabular}
    \caption{XX 班同学实践情况}
\end{table}

\subsubsection{社工}

此处介绍班级同学在学校和院系各类组织中担任职务的情况。

\begin{table}[H]
    \centering
    \zihao{-3}
    \begin{tabular}{||c|c||}
        \hline
        \hline
        姓名   & 岗位                       \\
        \hline
        XXX & 系团委 XX 组副书记         \\
            & 系课调委轮值主席           \\
        \hline
        XXX & 校 XX 协会负责人           \\
        \hline
        XXX & 系学生会 XX 部部长         \\
        \hline
        XXX & 党支部组织委员             \\
        \hline
        \hline
    \end{tabular}
    \caption{XX 班同学社工情况}
\end{table}
